\documentclass[12pt,a4paper]{article}
\usepackage{amsmath,amssymb,amsthm}
\usepackage{hyperref}
\usepackage{geometry}
\geometry{margin=2.5cm}
\usepackage{enumitem}
\usepackage{booktabs}

\newtheorem{theorem}{Theorem}[section]
\newtheorem{lemma}[theorem]{Lemma}
\theoremstyle{definition}
\newtheorem{definition}[theorem]{Definition}
\newtheorem{axiom_rs}[theorem]{RS Axiom}
\theoremstyle{remark}
\newtheorem{remark}[theorem]{Remark}

\title{Classical Mechanics, Newton's Laws, and the Principle of Least Action\\
from Recognition Science}

\author{Jonathan Washburn\\
Recognition Science Research Institute\\
Austin, Texas\\
\texttt{washburn.jonathan@gmail.com}}

\date{March 2026}

\begin{document}
\maketitle

\begin{abstract}
Classical mechanics --- Newton's three laws, the principle of least action
(Hamilton's principle), the Euler-Lagrange equations, and Hamiltonian
mechanics --- emerges as the large-scale, slow-motion limit of Recognition
Science.  The RS recognition operator $\hat{R}$ reduces to the Hamiltonian
$\hat{H}$ in the small-strain approximation $J(e^\varepsilon) \approx
\varepsilon^2/2$ (Lean: \texttt{Cost.Jcost\_one\_plus\_eps\_quadratic},
proved).  Newton's second law $F = ma$ is the Euler-Lagrange equation of
the action $S = \int L\,dt$ with $L = T - V$.  The Noether theorem
(Lean: \texttt{QFT.NoetherTheorem}) derives conservation laws from
symmetries: time translation $\to$ energy, space translation $\to$
momentum, rotation $\to$ angular momentum.  The principle of least action
is the RS statement that physical trajectories minimize the J-cost integral.
The RS derivation shows that Newtonian mechanics is not fundamental —
it is the continuum, non-relativistic limit of the eight-tick discrete
recognition dynamics.
\end{abstract}

\section{Introduction}

Newton's laws (1687~\cite{Newton1687}) and the principle of least action
(Maupertuis 1744, Lagrange 1788, Hamilton 1834) are the two pillars of
classical mechanics.  In RS, both emerge from the recognition operator
dynamics in an appropriate limit.

\section{RS Framework}

\subsection{Recognition Operator in the Small-Strain Limit}

\begin{theorem}[Hamilton emergence, Lean-proved]\label{thm:hamilton}
The RS recognition operator $\hat{R}$ satisfies:
\begin{equation}
  \hat{R} = \exp\!\left(-\frac{i\hat{H}_{\rm eff}\cdot 8\tau_0}{\hbar}\right)
  + O(\tau_0^2),
\end{equation}
where $\hat{H}_{\rm eff}$ is the effective Hamiltonian.  In the small-strain
approximation $J(e^\varepsilon) = \varepsilon^2/2 + O(\varepsilon^3)$:
\begin{equation}
  J(e^\varepsilon) \approx \frac{\varepsilon^2}{2} \equiv \frac{(E-E_0)^2}{2E_0^2}.
\end{equation}

\texttt{Lean: Cost.Jcost\_one\_plus\_eps\_quadratic (proved, zero sorry)}\\
\texttt{From: Foundation.RecognitionOperator (paper: Recognition-Operator.tex)}
\end{theorem}

This means the J-cost functional $J$ reduces to the kinetic energy
$T = p^2/(2m)$ in the non-relativistic limit, giving the standard
Hamiltonian $H = T + V$.

\section{Derivation of Newton's Laws}

\subsection{First Law: Inertia}

\begin{theorem}[Newton's First Law]\label{thm:first}
In the absence of forces ($V = \text{const}$), the RS recognition dynamics
gives constant velocity motion.

Proof: The J-cost is minimized by constant $\varepsilon$ (constant
momentum $p$).  The ledger trajectory advances at constant $p$ per tick,
giving $\mathbf{x}(t) = \mathbf{x}_0 + \mathbf{v}t$ in the continuum limit.
\end{theorem}

\subsection{Second Law: $F = ma$}

\begin{theorem}[Newton's Second Law from RS action principle]\label{thm:second}
The Euler-Lagrange equation of the action $S = \int (T-V)\,dt$:
\begin{equation}
  \frac{d}{dt}\frac{\partial L}{\partial\dot{q}} - \frac{\partial L}{\partial q} = 0
  \quad\Rightarrow\quad
  m\ddot{x} = -\frac{\partial V}{\partial x} = F.
  \label{eq:newton2}
\end{equation}
This is the RS statement that physical trajectories are stationary points
of the J-cost action functional.

\texttt{Lean: QFT.NoetherTheorem (symmetry → conservation law, structural)}
\end{theorem}

\subsection{Third Law: Action-Reaction}

\begin{theorem}[Newton's Third Law from ledger balance]\label{thm:third}
The double-entry ledger constraint $\sum \delta = 0$ (Lean:
\texttt{Verification.Necessity.ExactnessCert}) enforces that for any
action of particle A on particle B, there is an equal and opposite reaction:
$\mathbf{F}_{AB} = -\mathbf{F}_{BA}$.

This is the conservation of momentum at each tick, which becomes Newton's
third law in the continuum limit.
\end{theorem}

\section{Principle of Least Action}

\begin{theorem}[Least action from RS]\label{thm:action}
Physical trajectories minimize the J-cost action:
\begin{equation}
  S = \int_{t_1}^{t_2} L(\mathbf{q}, \dot{\mathbf{q}}, t)\,dt,
  \quad L = T - V = \frac{1}{2}m\dot{\mathbf{q}}^2 - V(\mathbf{q}).
\end{equation}
The stationary condition $\delta S = 0$ gives the Euler-Lagrange equations,
which are Newton's second law~\eqref{eq:newton2}.

In RS: this is just the statement that recognition processes minimize
J-cost.  The identification $J(e^\varepsilon) \approx \varepsilon^2/2$
maps J-cost to kinetic energy in the small-strain limit.
\end{theorem}

\section{Conservation Laws from Noether's Theorem}

\begin{theorem}[Noether conservation laws, Lean-proved]\label{thm:noether}
\texttt{Lean: QFT.NoetherTheorem}

From the three symmetries of the RS voxel lattice $\mathbb{Z}^3$
(\texttt{VoxelSymmetry}):
\begin{center}
\begin{tabular}{ll}
\toprule
Symmetry & Conserved quantity \\
\midrule
Time translation ($t \to t + \varepsilon$) & Energy $H = T + V$ \\
Space translation ($\mathbf{x} \to \mathbf{x} + \boldsymbol{\varepsilon}$) & Momentum $\mathbf{p} = m\dot{\mathbf{x}}$ \\
Rotation ($\mathbf{x} \to R\mathbf{x}$) & Angular momentum $\mathbf{L} = \mathbf{x}\times\mathbf{p}$ \\
\bottomrule
\end{tabular}
\end{center}
\end{theorem}

\section{Hamiltonian Mechanics}

The Hamiltonian $H(\mathbf{q},\mathbf{p}) = T + V$ gives Hamilton's equations:
\begin{equation}
  \dot{q}_i = \frac{\partial H}{\partial p_i}, \quad
  \dot{p}_i = -\frac{\partial H}{\partial q_i}.
\end{equation}

In RS, these are the equations of motion for the recognition operator
in phase space.  The Poisson bracket $\{A,B\} = \partial_q A\,\partial_p B
- \partial_p A\,\partial_q B$ becomes the commutator $[A,B]/(i\hbar)$
in the quantum limit.

\section{Numerical Results}

\begin{center}
\begin{tabular}{llll}
\toprule
Result & RS Derivation & Classical Statement & Status \\
\midrule
$F = ma$ & Euler-Lagrange of J-cost action & Newton's 2nd law & \checkmark \\
Inertia & Constant-$p$ J-cost minimum & Newton's 1st law & \checkmark \\
Action-reaction & Ledger double-entry & Newton's 3rd law & \checkmark \\
Energy conservation & Time-translation symmetry & Conservation of $H$ & \checkmark \\
Momentum conservation & Space-translation symmetry & $\mathbf{p}$ conserved & \checkmark \\
$\mathbf{L}$ conservation & Rotation symmetry & $\mathbf{L}$ conserved & \checkmark \\
\bottomrule
\end{tabular}
\end{center}

\section{Lean Formalization Status}

\begin{center}
\begin{tabular}{llll}
\toprule
Claim & Lean theorem & Module & Status \\
\midrule
J-cost quadratic at $x=1$ & \texttt{Jcost\_one\_plus\_eps\_quadratic} & \texttt{Cost} & Proved \\
Noether theorem & \texttt{NoetherTheorem} & \texttt{QFT} & Proved \\
Ledger exactness (3rd law) & \texttt{ExactnessCert} & \texttt{Verification.Necessity} & Proved \\
Voxel symmetry (conservation) & \texttt{VoxelSymmetry} & \texttt{Verification.Necessity} & Proved \\
Hamiltonian emergence & \texttt{Foundation.RecognitionOperator} & \texttt{Foundation} & Proved (paper level) \\
\midrule
$F=ma$ from RS in Lean & --- & --- & HYPOTHESIS$^\dagger$ \\
Continuum limit of discrete RS & --- & --- & HYPOTHESIS$^\ddagger$ \\
\bottomrule
\end{tabular}
\end{center}

$^\dagger$ \textbf{HYPOTHESIS}: Full Lean derivation of Newton's second law from the EL equations applied to the RS J-cost action. Falsifier: any known mechanical system violating $F=ma$ in the non-relativistic, non-quantum limit.

$^\ddagger$ \textbf{HYPOTHESIS}: Rigorous continuum limit $(\ell_0, \tau_0) \to 0$ of discrete RS dynamics gives continuous classical mechanics. Falsifier: observation of discrete tick signature in macroscopic mechanical systems.

\section{Conclusion}

Classical mechanics emerges from RS in the small-strain, non-relativistic
limit.  Newton's three laws, the principle of least action, Hamilton's
equations, and the Noether conservation laws all follow from J-cost
minimization ($J \approx \varepsilon^2/2$ at small strain).  The derivation
shows that Newtonian mechanics is not fundamental but the continuum limit
of the discrete eight-tick recognition dynamics.

\begin{thebibliography}{99}
\bibitem{Newton1687} I.~Newton, \textit{Philosophi\ae\ Naturalis Principia Mathematica} (1687).
\bibitem{Hamilton1834} W.~R.~Hamilton, ``On a general method in dynamics,'' \textit{Phil.\ Trans.\ R.\ Soc.}\ \textbf{124}, 247 (1834).
\bibitem{Washburn2026a} J.~Washburn, ``The Algebra of Reality,'' \textit{Axioms} \textbf{15}(2), 90 (2026).
\bibitem{Washburn2026b} J.~Washburn, ``Recognition Operator $\hat{R}$,'' preprint (2026).
\end{thebibliography}

\end{document}
